\documentclass[10pt,a4paper]{amsart}
\usepackage[margin=19mm]{geometry}
\usepackage{amsmath,amssymb,mathtools}
\usepackage[scr=rsfs,cal=euler]{mathalfa}
\usepackage{xfrac}
\usepackage[unicode,hidelinks,bookmarks=false]{hyperref}
\newcommand{\ie}{i.e.\ }
\newcommand{\cf}{cf.\ }
\newcommand{\resp}{respectively\ }
\newcommand{\Subsection}{\subsection}
\providecommand{\nobreakdash}{\nobreak\hbox{-}}
\setlength{\emergencystretch}{2em}
\multlinegap0pt
\usepackage{mathtools}
\usepackage{amssymb}
\usepackage{xcolor}
\newtheorem{lemma}{Lemma}
\title{Solvable groups defined by redundant cyclic presentations}
\author{Ihechukwu Chinyere}
\date{Source: arXiv:2609.08863v1 (CC BY 4.0)}
\begin{document}
\maketitle

\noindent\emph{Source and licence.} Solvable groups defined by redundant cyclic presentations. Version arXiv:2609.08863v1 (CC BY 4.0), distributed by arXiv under the Creative Commons Attribution 4.0 International licence, http://creativecommons.org/licenses/by/4.0/, as recorded in the arXiv licence metadata for this exact version and shown on the abstract page https://arxiv.org/abs/2609.08863v1. Full licence notice: \texttt{licenses/arxiv-2609.08863-CC-BY-4.0.txt}.\par\smallskip

\noindent\textit{Editorial scope.} Counted unit: the article's lemma lem:BSm-example (source0.tex lines 195--205). The article's complete original proof of it is delivered with it (source0.tex lines 207--259, one \texttt{proof} environment of 53 lines). No mathematics is written, repaired or re-derived: the statement and the proof are the article's own lines, in the article's own notation, and no retained line is substituted or re-flowed.  Retained uncounted as the source-local context the proof needs: lines 179--181.  Nothing was omitted from any retained span, so the package carries no omission record.  No retained line contains a citation, so the wrapper carries no bibliography and no bibliography entry.  No retained line contains a figure environment, an image-inclusion command or drawing code, so no image is delivered and the package declares no asset dependency.  Editorial formatting only, in addition: an AMS \texttt{amsart} preamble that restates the article's own declarations and macros used by the retained lines, namely the environments \texttt{lemma} on separate counters, together with the standard packages those lines need.  The retained environments print in this document as Lemma 1 in order of appearance, whereas the article prints the same items with the numbering of its own full document.  The complete native source of the article as distributed in the arXiv e-print for this exact version is bundled unchanged as \texttt{sources/209756/source0.tex}.
\begin{lemma}\label{lem:sigma-coords}
In the coordinates $(a,t)$, $\theta(t)=at$ and $\theta(a)=a^{-1}$.
\end{lemma}

\begin{lemma}\label{lem:BSm-example}
Let $m \in \mathbb{Z}$, and set $R_m := U_m(x_0,x_1)\,U_m(x_1,x_0)^{-1}$, where
\[
U_m(x_0,x_1) =
\begin{cases}
(x_1 x_0^{-1})^{m/2} x_1^{-1}, & m \text{ even}, \\[2pt]
x_0 (x_1 x_0^{-1})^{(m-1)/2} x_1, & m \text{ odd}.
\end{cases}
\]
Then $G_m =\langle x_0, x_1 \mid R_m \rangle$ is isomorphic to $BS(1,m) = \langle a, t \mid t^{-1} a t = a^m \rangle$.
\end{lemma}

\begin{proof}
Notice in particular, that $G_0\cong BS(1,0)\cong\mathbb{Z}$, $G_1\cong BS(1,1)\cong\mathbb{Z}^2$ and $G_{-1}\cong BS(1,-1),$ the Klein bottle group. As before, introduce the change of coordinates
\[
a = x_0 x_1^{-1}, \qquad t = x_1,
\]
so that $x_0 = at$, $x_1 = t$, and $x_0^{-1} = t^{-1}a^{-1}$. By Lemma~\ref{lem:sigma-coords}, the automorphism $\theta$ induced by the swap $x_0 \leftrightarrow x_1$ acts on these coordinates by
\[
\theta(a) = a^{-1}, \qquad \theta(t) = at.
\]
Set $P_m := U_m(at,t)$. Since $U_m(x_1,x_0)$ is obtained from $U_m(x_0,x_1)$ by applying $\theta$, we have $U_m(x_1,x_0) = \theta(P_m)$, and hence
\[
R_m = P_m\,\theta(P_m)^{-1}.
\]
We compute $P_m$ and $R_m$ separately according to the parity of $m$; in both cases, $x_1 x_0^{-1} = t(at)^{-1} = a^{-1}$.

\smallskip
\noindent\textbf{Case 1: $m = 2k$ even} (here $k \in \mathbb{Z}$ may be negative). Substituting,
\[
P_m = \bigl(t(at)^{-1}\bigr)^k t^{-1} = a^{-k} t^{-1}, \qquad
\theta(P_m) = a^{k}(at)^{-1} = a^{k} t^{-1} a^{-1}.
\]
Therefore
\[
R_m = P_m\,\theta(P_m)^{-1}
= a^{-k}t^{-1}\bigl(a^{k}t^{-1}a^{-1}\bigr)^{-1}
= a^{-k}\bigl(t^{-1}at\bigr)a^{-k}.
\]
Hence,  $R_m = 1$ holds if and only if $t^{-1}at = a^{2k} = a^m$, so
\[
G_m \cong \langle a, t \mid t^{-1}at = a^m \rangle = BS(1,m).
\]

\smallskip
\noindent\textbf{Case 2: $m = 2k+1$ odd} (here $k = (m-1)/2 \in \mathbb{Z}$ may be negative). Substituting,
\[
P_m = at\,a^{-k}t, \qquad
\theta(P_m) = a^{-1}(at)\,a^{k}(at) = t\,a^{k}(at) = t\,a^{k+1}t.
\]
Therefore
\[
R_m = P_m\,\theta(P_m)^{-1}
= at\,a^{-k}t\,\bigl(t\,a^{k+1}t\bigr)^{-1}
= at\,a^{-k}\bigl(t\,t^{-1}\bigr)a^{-(k+1)}t^{-1}
= at\,a^{-m}t^{-1}.
\]
Hence,  $R_m = 1$ holds if and only if $t a^m t^{-1} = a$, equivalently $t^{-1}at = a^m$, so again
\[
G_m \cong \langle a, t \mid t^{-1}at = a^m \rangle = BS(1,m).
\]

\smallskip
This proves $G_m \cong BS(1,m)$ in all cases. 
\end{proof}

\end{document}
